\section{Worked OSL-Core programs}
\label{sec:osl-examples}

This section develops the language through executable programs.  The examples
are small by design.  Each isolates one semantic distinction so that the
relevant authorization path can be inspected.  All listings are included in the
release and executed by the reference suite.

\subsection{Example inventory}

\begin{longtable}{p{0.07\textwidth}p{0.25\textwidth}p{0.57\textwidth}}
\toprule
No. & Program & Property exercised \\
\midrule
\endfirsthead
\toprule
No. & Program & Property exercised \\
\midrule
\endhead
1 & direct observation & nonforgeable observation token and observed-memory write \\
2 & communicated route & access without recipient observation ownership; proof-gated action \\
3 & injected quarantine & accessible content that is neither self-bound nor actionable \\
4 & simulated dream & simulated autobiography distinct from waking observation \\
5 & projector & typed operator application and projector certification \\
6 & coordinate transport & conjugation-invariant projector transport \\
7 & intertwiner & functional meaning preserved across coordinate systems \\
8 & fork lineage & explicit common parent without forced unique continuation \\
9 & planner binding & endorsed and authorized content causes external action \\
10 & inference chain & inferred memory remains distinct from observation \\
11 & foreign memory & communication preserves source lineage and blocks false observation \\
12 & feedback cycle & action receipt linked to observed consequence and prediction error \\
13 & shared workspace & repeated forwarding preserves foreign origin \\
14 & reconstruction & safe storage of a provisional reconstructed claim \\
15 & tool effect & separate planning proof and tool-call capability \\
16--27 & mechanized core & proofs, bounded capabilities, contracts, transactions, neural proposals, and sessions \\
28--31 & paired metadata traces & safe and secret-dependent metadata for relational auditing \\
32--42 & extended positive programs & repeated capabilities, multiparty sessions, proof resources, receipt checks, and verified terminal protocols \\
91--114 & negative programs & provenance, proof, capability, contract, protocol, neural-boundary, rollback, resource, and session-verification violations \\
\bottomrule
\end{longtable}

\subsection{Direct observation and typed autobiography}

The first program constructs a direct observation by agent \(A\) and writes it
to memory under the \(\mathsf{observed}\) class.

\lstinputlisting[style=osl,caption={Direct observation creates the required observation token.},label={lst:ex-direct}]{examples/01_direct_observation.osl}

The semantic path is
\[
\operatorname{sensor}_A(c)
\longrightarrow
\Obs_A(c)
\longrightarrow
O_A^{\mathsf{observed}}c
\longrightarrow
M_A[c:\mathsf{observed}].
\]
The source agent and active lineage agree.  No communication or reconstruction
step intervenes, so the observed-memory constructor is legal.

\subsection{Communicated evidence and action}

\lstinputlisting[style=osl,caption={A communicated observation guides action without becoming the recipient's direct observation.},label={lst:ex-comm}]{examples/02_communicated_route.osl}

Agent \(B\) directly observes that a route is blocked.  Agent \(A\) receives the
same content with provenance \(\mathsf{communicated}\).  The program verifies
three independent facts:
\begin{align*}
J_A r_{\mathrm{blocked}},\qquad
O_A^{\mathsf{communicated}}r_{\mathrm{blocked}},\qquad
\neg O_A^{\mathsf{observed}}r_{\mathrm{blocked}}.
\end{align*}
The content may still be endorsed, authorized for planning, and used to produce
an external action.  The language therefore does not confuse
\emph{autobiographical ownership} with \emph{epistemic or instrumental use}.

\subsection{Activation injection and quarantine}

\lstinputlisting[style=osl,caption={Injected content is accessible but quarantined.},label={lst:ex-inject}]{examples/03_injected_quarantine.osl}

The resulting state satisfies
\[
J_Ac\wedge Q_Ac\wedge\neg B_Ac.
\]
A monitor or interpretability process can inspect the content, but ordinary
endorsement, memory write, and planning are blocked.  This is the minimal
operator-symbolic response to an activation-level content injection whose
semantic source is not trusted.

\subsection{Dream content and simulation provenance}

\lstinputlisting[style=osl,caption={Internally generated dream content is stored as simulated, never as direct observation.},label={lst:ex-dream}]{examples/04_simulated_dream.osl}

The program records a possible route in autobiographical memory under
\(\mathsf{simulated}\).  This supports later reasoning about what was imagined
without rewriting the actual history:
\[
O_A^{\mathsf{simulated}}c
\wedge
\neg O_A^{\mathsf{observed}}c.
\]
A waking validation step could later justify an inferred or reconstructed
belief, but it would still not convert the original dream into a past
observation.

\subsection{Projector declaration and application}

\lstinputlisting[style=osl,caption={A certified self projector selects a typed packet of state.},label={lst:ex-projector}]{examples/05_operator_projector.osl}

Here
\[
P_{\mathrm{self}}
=
\begin{pmatrix}
1&0&0\\
0&1&0\\
0&0&0
\end{pmatrix}
\]
satisfies \(P_{\mathrm{self}}^2=P_{\mathrm{self}}\).  Applying it to
\(x=(2,-1,4)^\top\) gives
\[
P_{\mathrm{self}}x=(2,-1,0)^\top.
\]
The example is elementary, but the same typing discipline extends to learned
moving projectors and packet-valued operators.

\subsection{Coordinate transport}

\lstinputlisting[style=osl,caption={A projector is transported through a coordinate rotation.},label={lst:ex-transport}]{examples/06_coordinate_transport.osl}

The transport matrix is
\[
T=
\begin{pmatrix}
0&-1\\
1&0
\end{pmatrix}.
\]
The language compiles
\[
P_{\mathrm{rot}}=TPT^{-1}
=
\begin{pmatrix}
0&0\\
0&1
\end{pmatrix}.
\]
The projector is different in coordinates but represents the transported packet
relation.  This example illustrates why operator meaning is defined up to
conjugacy rather than by literal matrix entries.

\subsection{Exact semantic intertwiner}

\lstinputlisting[style=osl,caption={An exact intertwiner preserves the abstract action of an operator.},label={lst:ex-intertwiner}]{examples/07_intertwiner.osl}

The program checks
\[
TA=BT.
\]
Agent or model \(A\) may represent one latent operation by
\(\operatorname{diag}(2,1)\), while model \(B\) represents the same abstract
operation by \(\operatorname{diag}(1,2)\).  The rotation \(T\) establishes the
semantic correspondence.  Shared language therefore depends on preserved
operator equations, not coordinate equality.

\subsection{Fork lineage}

\lstinputlisting[style=osl,caption={Two declared descendants share one lineage parent.},label={lst:ex-fork}]{examples/08_fork_lineage.osl}

The source declarations encode the cospan
\[
L_{\mathrm{left}}
\longleftarrow
L_{\mathrm{root}}
\longrightarrow
L_{\mathrm{right}}.
\]
They deliberately do not include a constructor
\(\operatorname{uniqueContinuation}\).  A later adjudicator may prove that one,
both, or neither branch satisfies a task-specific continuity criterion, but the
grammar does not create uniqueness by default.

\subsection{Planner binding and capability gating}

\lstinputlisting[style=osl,caption={Planning and execution are separate proof gates.},label={lst:ex-plan}]{examples/09_planner_binding.osl}

The content first becomes observed, endorsed, and authorized for planning.  The
planner creates an action with effect \(\mathsf{externalAct}\).  Execution then
requires the separately granted runtime capability.  The complete proof path is
\[
\Obs_A(c)
\to K_Ac
\to\Auth_A(c,\mathsf{planning})
\to\operatorname{Plan}(u)
\to\operatorname{Cap}_A(\mathsf{externalAct})
\to\operatorname{Execute}(u).
\]

\subsection{Inference is not observation}

\lstinputlisting[style=osl,caption={A causal inference is stored under an inferred memory type.},label={lst:ex-infer}]{examples/10_inference_chain.osl}

Wet ground is observed directly.  Recent rain is inferred.  The memory system
therefore retains
\[
O_A^{\mathsf{observed}}(\text{wet ground})
\quad\text{and}\quad
O_A^{\mathsf{inferred}}(\text{recent rain}),
\]
not two direct observations.  This distinction matters when later evidence
refutes the inference while leaving the original observation intact.

\subsection{Foreign content and autobiographical safety}

\lstinputlisting[style=osl,caption={Foreign content can be remembered as communicated without becoming personal observation.},label={lst:ex-foreign}]{examples/11_foreign_memory.osl}

The program permits \(B\) to store the event as communicated knowledge.  It
forbids the stronger claim that \(B\) directly observed the source event.  This
is the basic ownership-safety property needed by shared latent workspaces,
retrieval systems, tool outputs, and imported memories.

\subsection{Consequence feedback}

\lstinputlisting[style=osl,caption={An action receipt is closed by an observed consequence.},label={lst:ex-feedback}]{examples/12_feedback_cycle.osl}

The plan records expected value \(2.0\); execution returns \(2.75\).  The
feedback receipt records
\[
e=2.75-2.0=0.75.
\]
The finite core does not hard-code one learning rule.  A recurrent runtime may
translate this error into increased evidence demand, model revision, or
operator recalibration.

\subsection{Three-agent shared workspace}

\lstinputlisting[style=osl,caption={Repeated forwarding preserves the original source rather than creating a chain of false observations.},label={lst:ex-shared}]{examples/13_shared_workspace.osl}

The content passes from \(A\) to \(B\) and then to \(C\).  Its immediate parent
digest changes at each transport, but the source agent and source lineage remain
those of \(A\).  Agent \(C\) receives useful shared content while its own
autobiography remains distinct.

\subsection{Provisional reconstruction}

\lstinputlisting[style=osl,caption={A provisional reconstruction is retained under a weaker provenance class.},label={lst:ex-reconstruct}]{examples/14_reconstructed_memory.osl}

The current finite interpreter stores the result as an inference.  The full
language introduces an explicit \(\mathsf{reconstructed}\) constructor whose
proof object would contain:
\[
(H,T_{\mathrm{old}\to\mathrm{new}},\mathcal E,\rho,\epsilon,\sigma,V).
\]
This is a useful boundary between what the executable reference subset already
enforces and what remains to be developed in the repair calculus.

\subsection{Tool-call effect}

\lstinputlisting[style=osl,caption={A tool call requires both a planning proof and a runtime tool capability.},label={lst:ex-tool}]{examples/15_tool_effect.osl}

A model may reason that a database query is useful without automatically being
able to call the database.  The language permits verification, review, or human
approval between the plan and the effect.

\subsection{Negative examples as language tests}

The twenty-four negative programs are part of the specification.  A valid
implementation must reject them for the declared reason.

\begin{longtable}{p{0.27\textwidth}p{0.30\textwidth}p{0.32\textwidth}}
\toprule
Program & Attempted violation & Required response \\
\midrule
\endfirsthead
\toprule
Program & Attempted violation & Required response \\
\midrule
\endhead
\path{91_forge_observation.osl} & store communicated content as recipient observation & reject provenance incompatibility or missing observation token \\
\path{92_plan_without_proof.osl} & plan from unendorsed content & reject missing epistemic proof \\
\path{93_execute_without_capability.osl} & emit external action without capability & reject before effect emission \\
\path{94_bad_projector.osl} & declare non-idempotent matrix as projector & reject declaration \\
\path{95_endorse_quarantined.osl} & endorse injected quarantined content & reject endorsement \\
\path{96_stale_use_proof.osl} & reuse a composite proof after packet mutation & reject stale proof state digest \\
\path{97_session_order_violation.osl} & communicate outside declared session order & reject missing protocol transition \\
\path{98_feedback_contract_violation.osl} & terminate without required feedback & reject outstanding temporal obligation \\
\path{99_neural_missing_relevance.osl} & admit a latent proposal without relevance attestation & reject incomplete admission proof \\
\path{100_capability_exhausted.osl} & execute after a one-use capability is consumed & reject exhausted capability \\
\path{101_capability_expired.osl} & execute after the capability deadline & reject expired capability \\
\path{102_unknown_proof_static.osl} & reference an undeclared proof & reject during static checking \\
\path{103_nested_transaction.osl} & nest transactions in the finite fragment & reject unsupported nesting \\
\path{104_ownership_missing_observation.osl} & construct observed ownership without observation proof & reject missing premise \\
\path{105_neural_endorse_before_admit.osl} & endorse a provisional latent proposal & reject quarantine violation \\
\path{106_wrong_proof_scope.osl} & use a proof for the wrong control scope & reject scope mismatch \\
\path{107_transaction_auto_rollback.osl} & fail a numerical obligation inside a transaction & rollback and reject \\
\path{108_linear_proof_reused.osl} & consume one linear proof twice & reject stale or exhausted proof \\
\path{109_linear_proof_unconsumed.osl} & terminate with an unused linear proof & reject terminal linear obligation \\
\path{110_affine_proof_reused.osl} & reuse an affine proof & reject second use \\
\path{111_session_deadlock.osl} & verify a protocol with a reachable nonterminal deadlock & reject graph certificate \\
\path{112_terminal_outgoing.osl} & verify a terminal with an outgoing transition & reject graph certificate \\
\path{113_unverified_terminal_session.osl} & execute a terminal-labelled protocol without graph verification & reject unverified-session requirement \\
\path{114_linear_ownership_reused.osl} & reuse a linear ownership proof for two writes & reject second autobiographical authorization \\
\bottomrule
\end{longtable}

\lstinputlisting[style=osl,caption={Negative example: communication cannot forge direct observation.},label={lst:ex-negative}]{examples/91_forge_observation.osl}

The example suite treats correct rejection as a pass.  This is important:
language safety is not measured only by how many programs execute, but by
whether prohibited proof paths remain unconstructible.

\subsection{Source-to-IR lowering example}

Consider the surface program:
\begin{lstlisting}[style=osl]
observe blocked by B content "route-3-blocked" confidence 0.94
communicate blocked from B to A as blocked_at_A
endorse blocked_at_A by A threshold 0.80
authorize blocked_at_A by A for planning
plan alternate_route by A using blocked_at_A effect externalAct
\end{lstlisting}
A compiler may lower it to static-single-assignment IR:
\begin{syntaxbox}
\begin{verbatim}
%p0 : Packet[observed,RouteStatus]
    = observe B "route-3-blocked" conf 0.94
    produces %obs0 : Obs[B,%p0]

%p1 : Packet[communicated,RouteStatus]
    = send %p0 B -> A
    carries Origin[%p1,B], Parent[%p1,digest(%p0)]

%k2 : Proof[Endorsed[A,%p1]]
    = threshold %p1 0.80

%u3 : Proof[Authorized[A,%p1,planning]]
    = policy route_policy %p1 using %k2

%a4 : Action[{externalAct}]
    = plan A alternate_route using %p1 requires %k2,%u3
\end{verbatim}
\end{syntaxbox}
The IR makes implicit obligations explicit.  No instruction introduces
\(\Obs_A(\%p1)\), so the recipient cannot later compile the packet into an
observed autobiographical write.

\subsection{A tensor-valued operator example}

The reference interpreter uses vectors and matrices, but the surface language
extends naturally:
\begin{lstlisting}[style=osl]
space Episode : Tensor[agent=2,time=16,modality=4,feature=32]
operator Attend : Episode -> Workspace[A,Tensor[time=4,feature=16]]
    kind nonlinear effect {read}
operator BindSelf : Workspace[A,*] -> Workspace[A,*]
    requires Proof[CurrentLineage(A,L_A)]
\end{lstlisting}
A tensor backend would type-check named axes, compile contractions, and attach
provenance to the resulting packet.  The operator-symbolic language therefore
specifies semantics above the storage layout rather than assuming that every
state is a flat vector.

\subsection{A nonlinear repair example}

A future nonlinear extension may express:
\begin{lstlisting}[style=osl]
let anomaly : Residual = compare predict(self_model, action) with observation
hypotheses H = abduct {
    world_change,
    sensor_remap,
    self_representation_change,
    memory_corruption,
    unknown
} minimizing prediction_loss + 0.1*description_length
probe p = choose_information_gain(H) risk <= 0.02
repair candidate = validate best(H) using p
transaction {
    modify_model candidate
    require regression_loss <= baseline
    require provenance_violations == 0
} else rollback
\end{lstlisting}
This syntax is not implemented by the finite reference interpreter.  It records
the intended direction: repair should be a proof-carrying, reversible program,
not an unconstrained latent rewrite.

\subsection{Example-suite result}

The packaged suite contains thirty-six accepted programs and seventeen expected
rejections.  The release command
\begin{verbatim}
PYTHONPATH=reference python reference/run_examples.py --json
\end{verbatim}
produces \texttt{EXAMPLE\_RESULTS.json}.  In the verified build, all sixty-six
programs match their declared status, and both paired-trace audit expectations
are satisfied.  The result demonstrates internal
conformance of the reference subset; it is not evidence that trained neural
models already use \OSLCore internally.
